\newif\ifuseseminar
\useseminartrue
\input{../../latex_common/header.tex.frag}

\title{Prova de Física Matemática II -- EDO}

\begin{document}
\begin{enumerate}
	\item Considere a equação diferencial linear de segunda ordem:
	      $$y'' + P_1(x)y' + P_0(x)y = 0.$$
	      Suponha que $x_0$ é um ponto singular regular. Podemos reescrever a equação na forma:
	      $$y'' + \frac{p(x)}{(x - x_0)}y' + \frac{q(x)}{(x - x_0)^2}y = 0,$$
	      onde $p(x) = (x - x_0)P_1(x)$ e $q(x) = (x - x_0)^2P_0(x)$ são funções analíticas em $x = x_0$. Com isso:
	      \begin{enumerate}
		      \item Determine o polinômio indicial e suas raízes a partir da forma de Frobenius.
		      \item Encontre a solução geral para a série de potências e derive a fórmula de recorrência.
		      \item Explique como o procedimento acima demonstra que sempre há pelo menos uma solução na forma de Frobenius.
	      \end{enumerate}
	\item Considere a equação de Bessel:
	      \[x^2y'' + xy' + (x^2 - \nu^2)y = 0.\]
	      \begin{enumerate}
		      \item Classifique o ponto $x_0 = 0$ e encontre a primeira solução usando o método de Frobenius.
		      \item Encontre a segunda solução no caso $\nu = 2$.
	      \end{enumerate}

	\item Considere a equação de Legendre:
	      $$
		      (1 - x^2)y'' - 2xy' + \lambda(\lambda + 1)y = 0,
	      $$
	      definida para $x \in [-1,1]$. Resolva a equação em $x_0 = 1$ usando o método
	      de Frobenius e encontre duas soluções linearmente independentes. Em seguida:
	      \begin{enumerate}
		      \item O que ocorre nos pontos singulares regulares se $\lambda$ não for um
		            número inteiro?
		      \item Quais são as condições sobre $\lambda$ para garantir que as soluções
		            sejam finitas em $x \in [-1,1]$?
	      \end{enumerate}
\end{enumerate}
\end{document}
